\clearpage
\section{Conclusion} \label{sec:conclusion}
  % Revisit the problem context
  The present study addresses the practical challenge of evaluating memory 
  disturbance effects on contemporary Raspberry Pi platforms. Specifically, 
  it focuses on implementing a Rowhammer-based proof-of-concept to investigate 
  whether repeated memory accesses can induce observable bit flips on the 
  Raspberry Pi 4 and 5. This work bridges the gap between theoretical 
  vulnerability assessments and experimental verification on modern embedded 
  ARM systems.
  \\
  % Degree of Objective Fulfillment
  Regarding the degree of objective fulfillment, the proof-of-concept was successfully 
  designed and implemented on both the Raspberry Pi 4 and 5 platforms. The implementation 
  incorporates two practical cache-bypass techniques - eviction-based strategies and cache 
  maintenance instructions, enabling uncached memory access for experimental hammering attempts.
  Furthermore, the \gls{dram} mapping functionality was fully integrated using the ported DRAMA 
  tool for ARM, allowing targeted memory row analysis on each test device. TRRespass 
  was successfully ported but did not generate effective access patterns, demonstrating 
  the current limitations of inducing Rowhammer bit flips on modern Raspberry Pi hardware.
  \\
  While all core modules of the proof-of-concept are operational and verified, inherent 
  hardware protections such as \gls{trr}, or page policy 
  behavior prevent reproducible bit flips, limiting the observable impact 
  of the implemented hammering techniques.
  \\
  This work provides a fully operational and documented Rowhammer proof-of-concept 
  specifically tailored to the Raspberry Pi 4 and 5 across multiple RAM configurations of 
  4~GB, 8~GB, and 16~GB. It establishes an ARM-adapted and publicly available combination 
  of the DRAMA and TRRespass tools, thereby facilitating further research on embedded 
  platforms. Moreover, the results confirm that, under standard operating conditions, 
  modern Raspberry Pi devices exhibit resilience against conventional software-induced 
  Rowhammer attacks, underscoring the effectiveness of contemporary embedded mitigation 
  mechanisms.
